\documentclass{article}

\usepackage{amsmath,amssymb}
\usepackage{xurl}
\usepackage[colorlinks=true,linkcolor=blue,citecolor=blue,urlcolor=blue]{hyperref}

% Shared commands for notes in docs/paper-gaps/.

\DeclareUnicodeCharacter{2080}{\ifmmode{}_{0}\else\textsubscript{0}\fi}
\DeclareUnicodeCharacter{2081}{\ifmmode{}_{1}\else\textsubscript{1}\fi}
\DeclareUnicodeCharacter{2082}{\ifmmode{}_{2}\else\textsubscript{2}\fi}
\DeclareUnicodeCharacter{2099}{\ifmmode{}_{n}\else\textsubscript{n}\fi}

\newcommand{\Lean}{\textsc{Lean}}
\ExplSyntaxOn
\NewDocumentCommand{\leanid}{m}{
  \ifmmode
    \textsf{\detokenize{#1}}
  \else
    \group_begin:
    \urlstyle{sf}
    \tl_set:Nn \l_tmpa_tl {#1}
    \tl_replace_all:Nnn \l_tmpa_tl {\_} {_}
    \exp_args:NV \path \l_tmpa_tl
    \group_end:
  \fi
}
\ExplSyntaxOff
\providecommand{\tr}{\operatorname{tr}}
\newcommand{\tnleanIssueBase}{https://github.com/LionSR/TNLean/issues/}
\newcommand{\tnleanPrBase}{https://github.com/LionSR/TNLean/pull/}
\newcommand{\tnleanDocBase}{https://sirui-lu.com/TNLean/docs/}
\newcommand{\ghissue}[1]{\href{\tnleanIssueBase#1}{GitHub issue \##1}}
\newcommand{\ghpr}[1]{\href{\tnleanPrBase#1}{PR \##1}}
\newcommand{\doclink}[1]{\href{\tnleanDocBase#1.html}{\path{#1}}}

% Dirac notation and the identity operator, as in blueprint/src/macros/common.tex.
% \providecommand keeps note-local definitions authoritative.
\usepackage{dsfont}
\providecommand{\ket}[1]{|#1\rangle}
\providecommand{\bra}[1]{\langle #1|}
\providecommand{\braket}[2]{\langle #1 | #2 \rangle}
\providecommand{\ketbra}[2]{|#1\rangle\!\langle #2|}
\providecommand{\Id}{\mathds{1}}
\providecommand{\one}{\mathds{1}}
\providecommand{\C}{\mathbb{C}}
\providecommand{\MN}[1]{M_{#1}(\C)}
\providecommand{\MD}{M_D(\C)}

% External referencing for paper sources.  The build helper writes sanitized
% auxiliary files under build/paper-aux/ and excludes duplicated source labels.
\usepackage{xr}

% Import a generated paper auxiliary file with a caller-supplied label prefix.
% The candidate paths support builds from docs/paper-gaps/, the repository root,
% and build/paper-aux/ itself.
\newcommand{\importpaperaux}[2]{%
  \IfFileExists{../../build/paper-aux/#2.aux}{%
    \externaldocument[#1]{../../build/paper-aux/#2}%
  }{%
    \IfFileExists{build/paper-aux/#2.aux}{%
      \externaldocument[#1]{build/paper-aux/#2}%
    }{%
      \IfFileExists{#2.aux}{%
        \externaldocument[#1]{#2}%
      }{}%
    }%
  }%
}

% General external reference macros
\newcommand{\paperref}[2]{\ref{#1-#2}}
\newcommand{\papereqref}[2]{(\ref{#1-#2})}

% CPSV16 source (1606.00608 / MPDO-22-12-17-2.tex) external document and shortcuts
\importpaperaux{cpsv16-}{1606.00608/MPDO-22-12-17-2-tnlean}
\newcommand{\cpsvref}[1]{\paperref{cpsv16}{#1}}
\newcommand{\cpsveqref}[1]{\papereqref{cpsv16}{#1}}

% Verdict marker: \gapnote{<kind>}{<status>}. Kind is the policy
% classification (clarification, local-correction, scope-restriction,
% unfaithful, false-source, open-gap); status is open, wip (elimination
% underway), resolved, or historical. Machine-read by the site index;
% prints a verdict line.
\newcommand{\gapnote}[2]{\par\noindent\textbf{Verdict:} #1 (#2).\par}


\title{One-Dimensional Scope of Schumacher--Werner Support-Algebra Covariance}
\author{}
\date{2026-09-03}

\begin{document}
\maketitle
\gapnote{scope-restriction}{open}

\section{The source identity}

Schumacher--Werner define a quantum cellular automaton on the cube lattice
$\mathbb Z^s$ with a finite neighbourhood scheme
\cite[lines~309--318]{SchumacherWerner2004QCA}. For their main structure
theorem they specialize to the cubic nearest-neighbour scheme, group the
$2^s$ sites of the unit cube $\Box$, and let $Q=\{-1,+1\}^s$ denote the
quadrant vectors \cite[lines~920--941]{SchumacherWerner2004QCA}. For each
$q\in Q$, they define the support algebra
\begin{align}
    \mathcal B_q
     & =\operatorname{Spp}\bigl(\gamma(\mathcal A(\Box)),
    \mathcal A(\Box+q)\bigr).
    \label{eq:sw04_support_covariance_scope_support}
\end{align}
Here $\gamma$ is the translation-commuting algebra homomorphism of their QCA
definition; invertibility is not assumed. Their translated
support calculation states that
\begin{align}
    \operatorname{Spp}(  \gamma(\mathcal A(\Box+q_1-q)),\mathcal A(\Box+q_1))
     & =\tau_{q_1-q}(\mathcal B_q).
    \label{eq:sw04_support_covariance_scope_source}
\end{align}
The definition and calculation occur at lines~1199--1206 and 1218--1226 of
Ref.~\cite{SchumacherWerner2004QCA}. They apply in arbitrary lattice
dimension with this cubic nearest-neighbour scheme.

\section{The formal specialization}

The present quasi-local observable algebra is a homogeneous chain on
$\mathbb Z$ with one fixed matrix size $d>0$. For a star-automorphism
$\omega$ satisfying
\begin{align}
    \omega(\mathcal A_\Lambda)
     & \subseteq\mathcal A_{\Lambda+\{-1,0,1\}},
    \label{eq:sw04_support_covariance_scope_propagation} \\
    \omega\tau_a
     & =\tau_a\omega
    \qquad(a\in\mathbb Z),
    \label{eq:sw04_support_covariance_scope_covariance}
\end{align}
the support algebras are formed from the even source pair
$E_x=\{2x,2x+1\}$ and its two neighbouring target pairs. Translation by two
sites preserves this pairing and gives
\begin{align}
    \tau_2(\widehat{\mathcal R}_y)
     & =\widehat{\mathcal R}_{y+2}.
    \label{eq:sw04_support_covariance_scope_formal}
\end{align}
This is the one-dimensional, nearest-neighbour, homogeneous-chain instance of
(\ref{eq:sw04_support_covariance_scope_source}). The fixed
on-site dimension agrees with the QCA setting used by Cirac et al. at
lines~2292--2298 of Ref.~\cite{Cirac2017MPU}; the ordered even and odd support
factors agree with equation \texttt{RR2x} of
Ref.~\cite{Gross2012IndexTheory}.

The formal statement assumes the propagation and translation covariance in
(\ref{eq:sw04_support_covariance_scope_propagation})--%
(\ref{eq:sw04_support_covariance_scope_covariance}), but it packages $\omega$
as a star-automorphism. Schumacher--Werner require only a translation-commuting
algebra homomorphism. Thus the formal theorem is additionally restricted to
the invertible case and does not formalize the arbitrary-dimensional cube
lattice.\footnote{Formal declaration:
    \leanid{SpinChain.PropagatesWithin.embeddedSupportAlgebra_map_quasiLocalTranslation_two}
    in \doclink{TNLean/QCA/SupportAlgebraCovariance}.}
The separate restriction from site-dependent GNVW cell sizes to a homogeneous
chain is recorded in
\path{docs/paper-gaps/gnvw12_site_dependent_adjacent_generation_scope.tex}.

\section{Elimination plan}

Removing the invertibility restriction first requires homomorphism-valued
versions of finite propagation, finite local restriction, translation
covariance, and the site-indexed support-algebra family. The compression proof
uses multiplication, addition, the star operation, and translation commutation;
it does not require an inverse for $\omega$.

Removing the dimensional restriction then requires a quasi-local observable
algebra indexed by $\mathbb Z^s$, finite-region translations by vectors in
$\mathbb Z^s$, and a propagation predicate for the cubic nearest-neighbour
scheme. For each translated source region and each target cube, the support algebra
must then be embedded canonically into the multidimensional quasi-local
algebra. Translation covariance and the same coefficient-compression
argument yield the vector-valued analogue of
(\ref{eq:sw04_support_covariance_scope_formal}), recovering
(\ref{eq:sw04_support_covariance_scope_source}). The present
theorem follows by specializing to $s=1$, $\mathcal N=\{-1,0,1\}$, and the
two-site blocked pairing.

Until both interfaces are available, the source-labelled formal statement
remains restricted to star-automorphisms of the one-dimensional
nearest-neighbour chain described above.

\bibliographystyle{alpha}
\bibliography{references}

\end{document}